\section{De las congruencias a Selmer: construcción y condiciones locales}
\label{sec:construccion-selmer}

La Nota~\ref{sec:congruencias-selmer} separó el denominador de un
proyector, una extensión residual y su imagen en cohomología divisible.
Su comparación final suponía un mapa cuyo núcleo había que controlar.
En esta nota construimos ese mapa en una situación precisa: una
representación irreducible de dimensión dos cuyos polinomios
característicos son congruentes con la suma de dos caracteres
residualmente distintos. Las entradas de sus matrices producirán
un cociclo de orden exacto; unos vectores propios locales medirán
la pérdida necesaria para obtener una clase de Selmer.

El caso aritmético con el primo $37$ realizará una clase de Selmer
estricta no nula. Su existencia utiliza un teorema clásico de Ribet,
que recibimos como entrada externa. El cálculo no identifica por sí
solo el ideal de congruencias de toda una familia theta: al final
precisaremos qué información adicional requiere esa comparación.

\fuente{A93, §9, pp.~61--63, para la construcción de extensiones y
la comparación por filtraciones; B76, §§5.1--5.2, pp.~54--57, para
las multiplicidades, el retículo torcido y las longitudes normalizadas.
El antecedente clásico es Ribet~\cite{ribet-congruencias},
Teorema~1.3, p.~152, y Proposición~2.1, pp.~154--155.
Las pruebas algebraicas siguientes se desarrollan aquí.
La existencia aritmética del apartado~\ref{subsec:rs-ribet}
se cita; no se añade una lectura ni una ejecución de Lean.}

\subsection{Los datos y la conclusión que buscamos}

Sean $K/\Q_\ell$ una extensión finita, $\mathcal O$ su anillo de
enteros, $\varpi$ un uniformizante y $k=\mathcal O/\varpi\mathcal O$.
Escribimos $R_n=\mathcal O/\varpi^n\mathcal O$, con $n\ge1$.
La cohomología será continua; los módulos de torsión llevan la
topología discreta. Fijamos un grupo profinito $G$, una representación
continua e irreducible sobre $K$
\[
 \rho:G\longrightarrow\operatorname{GL}_2(K)
\]
y caracteres continuos $\chi_1,\chi_2:G\to\mathcal O^\times$.
Suponemos que $\overline\chi_1\ne\overline\chi_2$ como caracteres
y que, para \emph{todo} $g\in G$,
\begin{equation}
 \begin{split}
 \operatorname{tr}\rho(g)&\equiv\chi_1(g)+\chi_2(g)\pmod{\varpi^n},\\
 \det\rho(g)&\equiv\chi_1(g)\chi_2(g)\pmod{\varpi^n}.
 \end{split}
 \label{eq:rs-congruencia}
\end{equation}
La traza y el determinante son integrales, como justificará el
Lema~\ref{lem:rs-reticulo}. La distinción residual exige un elemento
$g_0$ con $\chi_1(g_0)-\chi_2(g_0)\in\mathcal O^\times$;
no exige esa separación para cada elemento de $G$.

Pondremos
\begin{equation}
 \eta=\chi_1\chi_2^{-1},\qquad A=(K/\mathcal O)(\eta).
 \label{eq:rs-coeficientes}
\end{equation}
El torcido indica que $g$ actúa por multiplicación por $\eta(g)$.
Construiremos una inyección $R_n\hookrightarrow H^1(G,A)$ y,
con datos locales adicionales, un mapa a un grupo de Selmer con
núcleo de longitud calculable. La irreducibilidad se refiere a
$\rho$ en característica cero: su reducción puede ser una extensión
no escindida de caracteres.

\subsection{El núcleo del cambio de coeficientes}

Antes de construir una clase, conviene saber qué puede anularla.
La aplicación
\begin{equation}
 j_n:R_n(\eta)\longrightarrow A,\qquad
 x\longmapsto\varpi^{-n}x+\mathcal O
 \label{eq:rs-jn}
\end{equation}
identifica $R_n(\eta)$ con $A[\varpi^n]$. Aquí se utiliza el
carácter $\eta$ ya fijado con valores en $\mathcal O^\times$;
un carácter definido solo módulo $\varpi^n$ no define, sin más,
una acción sobre todo $K/\mathcal O$.

\begin{lema}[Obstrucción exacta y separación residual]
\label{lem:rs-coeficientes}
Para cualquier carácter continuo $\eta:G\to\mathcal O^\times$ se tiene
una sucesión exacta
\begin{equation}
 0\longrightarrow A^G/\varpi^n A^G
 \longrightarrow H^1(G,R_n(\eta))
 \xrightarrow{j_{n,*}}H^1(G,A)[\varpi^n]
 \longrightarrow0.
 \label{eq:rs-nucleo-coef}
\end{equation}
Si existe $g_0$ con $\eta(g_0)-1$ unidad, entonces $A^G=0$ y
$j_{n,*}$ es inyectivo. Si $\eta\ne1$ y
$s=\min_g v_\varpi(\eta(g)-1)$, con $v_\varpi(0)=+\infty$,
el núcleo tiene longitud $\min(n,s)$. Si $\eta=1$, el núcleo
es cero.
\end{lema}

\begin{proof}
La multiplicación por $\varpi^n$ en $A$ es sobreyectiva. La sucesión
larga de cohomología de
$0\to A[\varpi^n]\to A\xrightarrow{\varpi^n}A\to0$
da \eqref{eq:rs-nucleo-coef}. Si $\eta(g_0)-1$ es unidad,
todo invariante $a$ satisface $(\eta(g_0)-1)a=0$, de donde $a=0$.

Si $\eta\ne1$, el ideal generado por los $\eta(g)-1$ es
$(\varpi^s)$, para un entero finito $s\ge0$. Por tanto
$A^G=A[\varpi^s]\simeq\mathcal O/\varpi^s\mathcal O$ y su
cociente por $\varpi^n A^G$ tiene longitud $\min(n,s)$.
Si $\eta=1$, se tiene $A^G=A$, que es divisible, y ese cociente
es cero.
\end{proof}

La obstrucción es el cociente de invariantes de
\eqref{eq:rs-nucleo-coef}. La coincidencia de caracteres residuales
no basta para declarar que existe un núcleo: la acción trivial
muestra por qué hay que conservar el cociente completo.
En nuestra situación $\eta(g_0)-1$ es unidad, de modo que la
clase que construiremos no sufrirá la anulación observada en
el modelo de la Nota~\ref{sec:congruencias-selmer}.

\subsection{Un retículo estable por compacidad}

\begin{lema}[Integralidad de una representación compacta]
\label{lem:rs-reticulo}
La representación $\rho$ admite un retículo $G$-estable
$L\subset K^2$. En una base de $L$, todas sus matrices están en
$\operatorname{GL}_2(\mathcal O)$.
\end{lema}

\begin{proof}
Partimos de un retículo $L_0$. La imagen de $G$ es compacta;
sus entradas matriciales están acotadas. Existe $r\ge0$ tal que
$\rho(g)L_0\subseteq\varpi^{-r}L_0$ para todo $g$. Definimos
$L=\sum_{g\in G}\rho(g)L_0$, donde la suma significa el submódulo
generado por todos esos retículos. Entonces
$L_0\subseteq L\subseteq\varpi^{-r}L_0$.
Como $\mathcal O$ es noetheriano, $L$ es finitamente generado;
es libre de rango dos por ser un módulo sin torsión sobre un DVR.
Además $\rho(h)L=L$ para todo $h$, pues la multiplicación por $h$
permuta los índices de la suma. Las matrices y sus inversas son
integrales, lo que prueba la conclusión.
\end{proof}

Este primer retículo no tiene todavía por qué mostrar una extensión
no escindida. Lo modificaremos de una forma que preserve su
estabilidad y normalice uno de los bloques de las matrices.

\subsection{Separar las diagonales sin introducir denominadores}

\begin{lema}[Un elemento separador y dos trazas]
\label{lem:rs-separador}
Puede elegirse una base integral $e_1,e_2$ de $L$ en la que
$\rho(g_0)=\operatorname{diag}(\lambda_1,\lambda_2)$,
con $\lambda_i\equiv\chi_i(g_0)\pmod{\varpi^n}$ y
$\lambda_1-\lambda_2$ unidad. Si escribimos
\begin{equation}
 \rho(g)=\begin{pmatrix}a(g)&b(g)\\c(g)&d(g)\end{pmatrix},
 \label{eq:rs-bloques}
\end{equation}
entonces $a(g)\equiv\chi_1(g)$ y
$d(g)\equiv\chi_2(g)\pmod{\varpi^n}$ para todo $g$.
\end{lema}

\begin{proof}
El polinomio característico de $\rho(g_0)$ tiene dos raíces
simples módulo $\varpi$, las reducciones de $\chi_i(g_0)$.
La forma de raíces simples del lema de Hensel produce raíces
$\lambda_i\in\mathcal O^\times$: sucesivamente se corrige una
raíz módulo $\varpi^r$ mediante la derivada, que es unidad,
y la completitud permite tomar el límite.
Al evaluar el polinomio en $\chi_i(g_0)$ obtenemos
\[
 (\chi_i(g_0)-\lambda_i)(\chi_i(g_0)-\lambda_j)
 \in\varpi^n\mathcal O\qquad(j\ne i).
\]
El segundo factor es unidad. Así se obtiene la precisión
$\lambda_i\equiv\chi_i(g_0)\pmod{\varpi^n}$.

Los operadores
\[
 P_1=\frac{\rho(g_0)-\lambda_2 I}{\lambda_1-\lambda_2},
 \qquad P_2=I-P_1
\]
son idempotentes complementarios y preservan $L$, porque el
denominador es unidad. La descomposición $L=P_1L\oplus P_2L$
tiene dos sumandos libres de rango uno. Elegimos un generador
en cada uno.

Aplicamos \eqref{eq:rs-congruencia} a $g$ y a $g_0g$. Se obtiene,
módulo $\varpi^n$,
\[
 \begin{split}
 a(g)+d(g)&=\chi_1(g)+\chi_2(g),\\
 \lambda_1a(g)+\lambda_2d(g)
 &=\chi_1(g_0)\chi_1(g)+\chi_2(g_0)\chi_2(g).
 \end{split}
\]
El determinante del sistema es una unidad, y sus soluciones son
las diagonales afirmadas. No hace falta que $g$ y $g_0$ conmuten:
solo usamos $\rho(g_0g)=\rho(g_0)\rho(g)$.
\end{proof}

Las raíces $\lambda_i$ y los valores $\chi_i(g_0)$ son
\emph{congruentes}, no necesariamente iguales en $K$.
El argumento no divide por $2$. También sirve en característica
residual $2$ si el cuerpo residual permite dos caracteres distintos;
por ejemplo, esto puede ocurrir sobre $\mathbf F_4$. Sobre
$\mathbf F_2$ todos los caracteres con valores en
$\mathbf F_2^\times$ coinciden, por lo que nuestra hipótesis no se
cumple allí.

\subsection{Normalizar el bloque que producirá la extensión}

\begin{proposicion}[Un retículo con bloque superior primitivo]
\label{prop:rs-normalizar}
Existe un retículo estable $\Lambda$ con una base en la que
\eqref{eq:rs-bloques} satisface
\begin{equation}
 \begin{gathered}
 a(g)\equiv\chi_1(g),\quad d(g)\equiv\chi_2(g)
       \pmod{\varpi^n},\qquad c(g)\in\varpi^n\mathcal O,\\
 (b(g):g\in G)=\mathcal O,\qquad b(g_0)=c(g_0)=0.
 \end{gathered}
 \label{eq:rs-normalizado}
\end{equation}
\end{proposicion}

\begin{proof}
En la base del lema anterior, sean $B=(b(g):g\in G)$ y
$C=(c(g):g\in G)$. La identidad
$a(gh)=a(g)a(h)+b(g)c(h)$ y la multiplicatividad de $\chi_1$
dan $b(g)c(h)\in\varpi^n\mathcal O$ para todo par $g,h$.
Por tanto $BC\subseteq\varpi^n\mathcal O$.
Los dos ideales son no nulos: $B=0$ haría estable la recta $Ke_2$;
$C=0$ haría estable $Ke_1$, contra la irreducibilidad.

Escribimos $B=\varpi^u\mathcal O$, $C=\varpi^v\mathcal O$,
con $u,v\ge0$ y $u+v\ge n$. Tomamos
$\Lambda=\mathcal O(\varpi^u e_1)+\mathcal O e_2$.
En esa base las cuatro entradas pasan a ser
\[
 a(g),\quad\varpi^{-u}b(g),\quad\varpi^u c(g),\quad d(g).
\]
Son integrales; el bloque inferior pertenece a
$\varpi^{u+v}\mathcal O\subseteq\varpi^n\mathcal O$ y el
superior genera $\mathcal O$. Todas las matrices, incluidas las
correspondientes a $g^{-1}$, son integrales; $\Lambda$ es estable.
Las diagonales no cambian y $g_0$ sigue diagonal. Omitimos las
primas para obtener \eqref{eq:rs-normalizado}.
\end{proof}

La irreducibilidad sobre $K$ impide que desaparezca uno de los dos
bloques antes de normalizarlo. Después de cambiar el retículo, el
bloque inferior desaparece módulo $\varpi^n$ y el superior conserva
al menos una unidad. Esa asimetría es la que hará visible la clase.

\subsection{La inyección construida y su anulador exacto}

\begin{teorema}[De una congruencia separada a una clase de longitud $n$]
\label{thm:rs-inyeccion}
Bajo \eqref{eq:rs-congruencia} y las hipótesis del primer apartado,
el retículo de la Proposición~\ref{prop:rs-normalizar} proporciona
el cociclo
\begin{equation}
 z_n(g)=b(g)\chi_2(g)^{-1}\pmod{\varpi^n}
       \quad\text{en }R_n(\eta).
 \label{eq:rs-cociclo}
\end{equation}
Las aplicaciones
\begin{equation}
 \begin{aligned}
 R_n&\longrightarrow H^1(G,R_n(\eta)),& u&\longmapsto u[z_n],\\
 R_n&\longrightarrow H^1(G,A),& u&\longmapsto u j_{n,*}[z_n]
 \end{aligned}
 \label{eq:rs-inyecciones}
\end{equation}
son inyectivas. Ambas clases tienen anulador exactamente
$(\varpi^n)$ como módulos sobre $\mathcal O$.
\end{teorema}

\begin{proof}
La entrada superior derecha de un producto da
$b(gh)=a(g)b(h)+b(g)d(h)$. Dividiendo por $\chi_2(gh)$ y
reduciendo módulo $\varpi^n$ obtenemos
\[
 z_n(gh)=z_n(g)+\eta(g)z_n(h).
\]
La continuidad procede de la de las entradas y de la reducción
al módulo finito. Es, por tanto, un cociclo continuo.

Supongamos $u[z_n]=0$. Existe $t\in R_n$ tal que
$u z_n(g)=(\eta(g)-1)t$ para todo $g$.
Al evaluar en $g_0$, el lado izquierdo es cero y
$\eta(g_0)-1$ es unidad; por tanto $t=0$.
Como las entradas $b(g)$ generan $\mathcal O$, alguna $b(g_1)$
es unidad. La igualdad en $g_1$ implica $u=0$ en $R_n$.
Esto prueba la primera inyección y el anulador exacto.
La segunda sigue del Lema~\ref{lem:rs-coeficientes}, pues
$\overline\chi_1\ne\overline\chi_2$.
\end{proof}

La reducción del retículo encaja en
$0\to R_n(\chi_1)\to\Lambda/\varpi^n\Lambda
\to R_n(\chi_2)\to0$. El cociclo de esta extensión toma valores
en $\operatorname{Hom}(R_n(\chi_2),R_n(\chi_1))=R_n(\eta)$.
El teorema prueba más que la no escisión módulo $\varpi$:
ningún múltiplo de valoración menor que $n$ anula su clase.
La aplicación depende de las elecciones de base; su existencia
y la longitud de su imagen son las conclusiones necesarias.

\subsection{Los datos locales y la pérdida que producen}

Sean $\iota_w:D_w\to G$ una familia finita de homomorfismos
continuos. Escribimos $\operatorname{res}_w$ para la restricción
por $\iota_w$, y definimos el grupo con condición local estricta
\begin{equation}
 S_{\mathrm{estr}}(G,A;D_w)=
 \ker\!\left(H^1(G,A)\longrightarrow\prod_w H^1(D_w,A)\right).
 \label{eq:rs-selmer-estricto}
\end{equation}
En una aplicación aritmética $G$ será $G_{E,S_\ell}$ y
$D_w=G_{E_w}$, para $w\mid\ell$; las acciones locales se
obtienen por los homomorfismos naturales. Como en la
Nota~\ref{sec:congruencias-selmer}, la notación $G_{E,S}$
permite también ramificación en los lugares infinitos, que no
se enumeran en $S$.

\begin{proposicion}[Un vector propio local controla el núcleo]
\label{prop:rs-local}
Para cada $w$, supongamos dado un vector primitivo
$v_w=(x_w,y_w)\in R_n^2$ tal que
\begin{equation}
 \rho(\iota_w(d))v_w=\chi_2(\iota_w(d))v_w
       \quad\text{en }R_n^2\quad(d\in D_w).
 \label{eq:rs-vector-local}
\end{equation}
Primitivo significa que alguna coordenada es unidad. Sea
$t_w=v_\varpi(y_w)$ si $y_w\ne0$, con $0\le t_w<n$,
y $t_w=n$ si $y_w=0$. Pongamos $t=\max_w t_w$;
para una familia vacía tomamos $t=0$. Si
$\gamma_n=j_{n,*}[z_n]$, la aplicación
\begin{equation}
 R_n\longrightarrow S_{\mathrm{estr}}(G,A;D_w),\qquad
 u\longmapsto u\varpi^t\gamma_n
 \label{eq:rs-mapa-selmer}
\end{equation}
está bien definida. Su núcleo es
$\varpi^{n-t}\mathcal O/\varpi^n\mathcal O$, de longitud $t$,
y su imagen tiene longitud $n-t$. Si el grupo de Selmer tiene
longitud finita, se obtiene
\begin{equation}
 n\le t+\length_{\mathcal O}S_{\mathrm{estr}}(G,A;D_w).
 \label{eq:rs-cota-local}
\end{equation}
\end{proposicion}

\begin{proof}
La primera coordenada de \eqref{eq:rs-vector-local}, usando
$a\equiv\chi_1$, da
$(\chi_1(d)-\chi_2(d))x_w+b(d)y_w=0$ en $R_n$;
hemos omitido $\iota_w$ en la notación. Al dividir por
$\chi_2(d)$ resulta
\begin{equation}
 y_w z_n(d)=(1-\eta(d))x_w=\delta(-x_w)(d).
 \label{eq:rs-coborde-local}
\end{equation}
Si $y_w=\varpi^{t_w}u_w$ con $u_w$ unidad, esta igualdad
anula $\varpi^{t_w}\operatorname{res}_w[z_n]$.
Si $y_w=0$, la multiplicación por $\varpi^n$ la anula de todas
formas. Así $\varpi^t\operatorname{res}_w[z_n]=0$ para cada $w$.
El cambio de coeficientes respeta las restricciones, de modo que
$\varpi^t\gamma_n$ pertenece al grupo indicado.

Por el Teorema~\ref{thm:rs-inyeccion}, $\gamma_n$ tiene anulador
$(\varpi^n)$. Por ello
$u\varpi^t\gamma_n=0$ equivale a
$u\in\varpi^{n-t}R_n$. Las longitudes del núcleo y de la imagen
se calculan en el módulo cíclico $R_n$. La inclusión de la imagen
en Selmer prueba \eqref{eq:rs-cota-local}.
\end{proof}

Si todos los $y_w$ son unidades, las extensiones locales se
escinden y la aplicación es inyectiva. Una coordenada divisible
por $\varpi$ mide una pérdida posible: el argumento construye
un mapa con ese núcleo, sin afirmar que sea el mejor mapa ni que
$t$ sea la mínima pérdida posible. Con $t=n$ solo obtenemos
el mapa cero. Los vectores se piden módulo $\varpi^n$; no hace
falta que se levanten a rectas propias sobre $K$.

Esta proposición prueba la condición local cero. Una condición
finita de Bloch--Kato puede ser más amplia; sustituirla por cero
sin justificarlo cambiaría el problema. Para pasar a otra
condición local hay que identificar el subgrupo correspondiente.

\subsection{Cómo verificar las congruencias mediante Frobenius}

En este apartado y en el ejemplo tomamos Frobenius aritmético;
por tanto $\chi_{\mathrm{cyc}}(\operatorname{Frob}_v)=Nv$,
el cardinal del cuerpo residual, cuando $v\nmid\ell$.
Para $E=\Q$ y un primo $q\ne\ell$, este valor es $q$.

\begin{proposicion}[Criterio de Frobenius]
\label{prop:rs-frobenius}
Supongamos $G=G_{E,S}$, donde $E$ es un cuerpo de números y $S$
un conjunto finito de lugares que contiene los situados sobre
$\ell$. Sean $\rho$ y $\chi_i$ representaciones de este grupo.
Si, fuera de un conjunto finito de lugares $v$, sus polinomios
característicos en Frobenius satisfacen
\begin{equation}
 \det(X-\rho(\operatorname{Frob}_v))\equiv
 (X-\chi_1(\operatorname{Frob}_v))
 (X-\chi_2(\operatorname{Frob}_v))\pmod{\varpi^n},
 \label{eq:rs-frobenius}
\end{equation}
entonces \eqref{eq:rs-congruencia} se cumple para todo $g\in G$.
\end{proposicion}

\begin{proof}
Elegimos un retículo estable. Las reducciones de $\rho$ y de los
dos caracteres módulo $\varpi^n$ tienen imagen finita y factorizan
por un mismo cociente finito de $G$. La traza y el determinante
son funciones de sus clases de conjugación. El teorema de
Chebotarev asegura que cada clase del cociente aparece como
Frobenius en lugares fuera de cualquier conjunto finito prefijado.
La igualdad en esas clases da la igualdad en todo el cociente.
\end{proof}

Se recibe aquí Chebotarev en su formulación usual para extensiones
finitas de cuerpos de números; véase la exposición de cuerpos
globales de Milne~\cite{milne-ant}, Teorema~8.31, p.~147. El argumento requiere congruencias para todos los lugares
salvo un conjunto finito. Comprobar una lista numérica de primos
no proporciona esta hipótesis.

Por ejemplo, cuando se dispone de una representación asociada a
un sistema de Hecke con polinomio
$X^2-a_qX+q^{k-1}$ y de congruencias
$a_q\equiv1+q^{k-1}\pmod{\varpi^n}$ en los lugares requeridos,
los caracteres naturales son $\chi_1=1$ y
$\chi_2=\chi_{\mathrm{cyc}}^{k-1}$. El cociclo toma entonces valores
en el torcido negativo $1-k$. Esta descripción presupone la
existencia e irreducibilidad de la representación y no establece
su escisión local en $\ell$.

\subsection{La entrada aritmética de Ribet}
\label{subsec:rs-ribet}

Para disponer de un ejemplo aritmético completo utilizamos el
siguiente resultado clásico, con sus condiciones locales incluidas.
Sea $p$ un primo impar y sea $k$ par, con $2\le k\le p-3$.
Escribimos $\omega:G_\Q\to\mathbf F_p^\times$ para el carácter
ciclotómico residual. Si $p$ divide el numerador de $B_k$,
el Teorema~1.3 de Ribet~\cite{ribet-congruencias}, p.~152,
proporciona un cuerpo finito $F\supseteq\mathbf F_p$ y una
representación continua
\begin{equation}
 \overline\rho:G_\Q\longrightarrow\operatorname{GL}_2(F),
 \qquad
 \overline\rho(g)=
 \begin{pmatrix}1&\beta(g)\\0&\omega(g)^{k-1}\end{pmatrix},
 \label{eq:rs-ribet}
\end{equation}
no ramificada fuera de $p$, no escindida globalmente y
escindida al restringirla a un grupo de descomposición en $p$.
En el enunciado original las dos últimas propiedades se expresan
mediante el orden de las imágenes: el global es divisible por $p$
y el local en $p$ es primo con $p$.

La fórmula
\begin{equation}
 z(g)=\beta(g)\omega(g)^{1-k}
 \label{eq:rs-ribet-cociclo}
\end{equation}
da una clase no nula en
$H^1(G_{\Q,\{p\}},F(\omega^{1-k}))$.
La identidad de cociclo se obtiene multiplicando matrices,
como en \eqref{eq:rs-cociclo}. Una clase cero permitiría cambiar
el levantamiento de la segunda coordenada para diagonalizar la
representación; esto contradice la no escisión global.
La escisión local implica que su restricción en $p$ es un coborde.

La existencia de \eqref{eq:rs-ribet} contiene la parte profunda del
ejemplo. Ribet explica en p.~153 que su construcción local utiliza
formas de peso $2$ y nivel $p$. Allí distingue este camino del
que parte de formas de peso $k$ y nivel $1$: las congruencias de
ese segundo camino no proporcionan por sí solas la propiedad
local requerida. Por ello no deducimos la escisión en $p$ de la
sola congruencia $a_q\equiv1+q^{k-1}$.

Podemos además obtener una clase sobre $\mathbf F_p$.
Elegimos una base de $F$ sobre $\mathbf F_p$ y sus aplicaciones
coordenadas $\lambda_i:F\to\mathbf F_p$.
Como la acción es por escalares de $\mathbf F_p$, todas son
equivariantes y la cohomología se descompone en esas coordenadas.
Si todas las clases $[\lambda_i z]$ fueran cero, $[z]$ sería cero.
Alguna coordenada $\lambda$ produce, por tanto, una clase no nula
y conserva la condición de coborde en $p$ y la no ramificación
fuera de $p$.

\subsection{El ejemplo aritmético con el primo 37}

Definimos los números de Bernoulli mediante $B_0=1$ y
\begin{equation}
 \sum_{j=0}^{m}\binom{m+1}{j}B_j=0\quad(m\ge1).
 \label{eq:rs-bernoulli}
\end{equation}
Así $B_1=-1/2$. La recurrencia, calculada en aritmética racional
exacta, da
\begin{equation}
 B_{32}=-\frac{7709321041217}{510}
       =-\frac{37\cdot683\cdot305065927}{510},
 \qquad v_{37}(B_{32})=1.
 \label{eq:rs-B32}
\end{equation}
El cociente del numerador por $37$ es
$-208360028141\equiv3\pmod{37}$ y $510\equiv29\pmod{37}$;
esto verifica que la divisibilidad es exactamente una vez.

\begin{proposicion}[Una clase estricta no nula]
\label{prop:rs-37}
Para $A=\Q_{37}/\Z_{37}(-31)$ existe una inyección
\begin{equation}
 \mathbf F_{37}\lhook\joinrel\longrightarrow
 S_{37}(\Q,31):=
 \ker\!\left(H^1(G_{\Q,\{37\}},A)
                  \longrightarrow H^1(G_{\Q_{37}},A)\right).
 \label{eq:rs-37-selmer}
\end{equation}
Su imagen es localmente trivial en todos los lugares; en particular
se inyecta en $\operatorname{Sha}^1(\Q,A)$, el núcleo de la restricción de
$H^1(G_\Q,A)$ al producto de todas las cohomologías locales.
\end{proposicion}

\begin{proof}
Aplicamos la entrada de Ribet con $p=37$ y $k=32$.
Se cumplen $2\le32\le34$ y la divisibilidad de
\eqref{eq:rs-B32}. El apartado anterior proporciona una clase no
nula $[z_0]=[\lambda z]$ con coeficientes
$\mathbf F_{37}(\omega^{-31})$, no ramificada fuera de $37$
y localmente cero en $37$.

La acción en $A$ es la del carácter
$\chi_{\mathrm{cyc}}^{-31}$. En un Frobenius aritmético en $2$,
su reducción vale
$2^{-31}=2^5=32\pmod{37}$, usando que
$-31\equiv5\pmod{36}$. Por tanto
$\chi_{\mathrm{cyc}}^{-31}(\operatorname{Frob}_2)-1$ es una
unidad de $\Z_{37}$. El Lema~\ref{lem:rs-coeficientes}, con $n=1$,
muestra que la imagen
$\gamma=j_{1,*}[z_0]$ en $H^1(G_{\Q,\{37\}},A)$ es no nula.
La aplicación explícita $a\mapsto a\gamma$ es inyectiva,
pues su dominio es el cuerpo $\mathbf F_{37}$. Su restricción
en $37$ es cero por functorialidad.

Para un primo $q\ne37$, la clase es no ramificada y $A$ también.
La cohomología no ramificada se calcula por
\begin{equation}
 H^1_{\mathrm{ur}}(\Q_q,A)
       =A/(q^{-31}-1)A=0.
 \label{eq:rs-local-q}
\end{equation}
En efecto, el grupo no ramificado está generado topológicamente
por Frobenius, lo que da el cociente; el escalar
$q^{-31}-1$ es un elemento no nulo de $\Z_{37}$ y la
multiplicación por él es sobreyectiva en el módulo divisible $A$.
En el lugar real, la cohomología de grado uno es anulada por el
orden $2$ del grupo de Galois; como $2$ es invertible en $A$,
también es cero. La inflación de $G_{\Q,\{37\}}$ a $G_\Q$ es
inyectiva en grado uno, por inflación--restricción, ya que la
acción en $A$ factoriza por ese cociente. Quedan comprobadas
todas las condiciones locales de $\operatorname{Sha}^1$.
\end{proof}

La conclusión es la existencia de un submódulo de longitud uno.
No hemos demostrado aquí la finitud de todo $S_{37}(\Q,31)$ ni
que su longitud sea exactamente uno. El valor
$v_{37}(B_{32})=1$ entra en la elección del ejemplo; una igualdad
con la longitud completa exigiría otro teorema aritmético.

También se obtiene una clase sobre $\Q(\zeta_{37})$.
La restricción seguida de corestricción multiplica por
$[\Q(\zeta_{37}):\Q]=36$, una unidad en $\Z_{37}$.
La restricción es, por ello, inyectiva y conserva la trivialidad
local. Este argumento no afirma que cualquier clase residual
sobre el cuerpo ciclotómico sobreviva al cambio de coeficientes:
restringimos una clase que ya hemos construido en el módulo
divisible.

\subsection{Qué ideal mide cada construcción}

La longitud $n$ de \eqref{eq:rs-inyecciones} pertenece a la
congruencia \eqref{eq:rs-congruencia} para los dos caracteres
fijados. En la prueba apareció además el ideal $BC$ de los
bloques. Se tiene $BC\subseteq(\varpi^n)$, pero puede haber
inclusión estricta; no hemos probado que $n=v_\varpi(BC)$.
Y ninguno de los dos se identifica automáticamente con el
ideal $\mathfrak c_\pi$ del proyector de Hecke estudiado en
las Notas~\ref{sec:cohomologia} y~\ref{sec:congruencias-selmer}.

La diferencia puede verse en un orden de Hecke abstracto
$\mathbb T\subseteq\mathcal O^X$, con proyecciones de coordenadas
$\lambda_\pi$. Consideramos el caso de un único sistema retenido,
$Z=\{\pi\}$; para un conjunto $Z$ mayor la condición de
anulación del proyector solo afecta al complemento de $Z$.
Para $\sigma\ne\pi$ definamos
\begin{equation}
 \begin{split}
 I_{\pi,\sigma}
  &=(\lambda_\pi(T)-\lambda_\sigma(T):T\in\mathbb T),\\
 \mathfrak c_\pi
  &=\{\lambda_\pi(T):T\in\mathbb T,
                      \ \lambda_\tau(T)=0\text{ para }\tau\ne\pi\}.
 \end{split}
 \label{eq:rs-ideales-hecke}
\end{equation}
Para cada $\sigma\ne\pi$ se cumple
\begin{equation}
 \mathfrak c_\pi\subseteq I_{\pi,\sigma},
 \qquad v_\varpi(\mathfrak c_\pi)\ge v_\varpi(I_{\pi,\sigma})
 \label{eq:rs-direccion-ideal}
\end{equation}
cuando los ideales son no nulos. Basta tomar un operador que
anule todas las otras coordenadas: su diferencia con la
coordenada $\sigma$ es precisamente su valor en $\pi$.
La dirección de la desigualdad importa. Construir una clase
para una sola pareja no da una cota superior para todo el
ideal del proyector. El Ejercicio~\ref{ej:rs-tres} exhibe una
separación arbitrariamente grande.

A93, p.~62, explica la construcción de extensiones a partir de
bloques. En rango dos,
normalizar $B=\mathcal O$ hace visible una copia de $R_n$.
Para una familia con más constituyentes, el módulo de bloques
puede tener más generadores y la comparación debe conservarlos.

\subsection{Filtraciones, multiplicidades y errores finitos}

El formato que necesita la aplicación general puede expresarse
mediante un lema de longitudes. Su contenido es elemental;
la existencia de los mapas es la entrada aritmética que debe
construirse en cada situación.

\begin{proposicion}[Contabilidad de una comparación filtrada]
\label{prop:rs-filtracion}
Sea $D$ un módulo de longitud finita con una filtración finita.
Supongamos que sus cocientes se dividen en dos tipos. Los del
primer tipo, $Q_i$, admiten mapas
\begin{equation}
 Q_i\longrightarrow\bigoplus_j S_j^{\oplus m_{ij}}
 \quad\text{con núcleo de longitud a lo sumo }\kappa_i,
 \label{eq:rs-graduados}
\end{equation}
donde los $S_j$ tienen longitud finita. Los del segundo tipo,
$E_h$, están anulados por $\varpi^{a_h}$ y admiten a lo sumo
$b_h$ generadores. Entonces
\begin{equation}
 \length_{\mathcal O}D
 \le\sum_i\kappa_i+
       \sum_{i,j}m_{ij}\length_{\mathcal O}S_j+
       \sum_h a_h b_h.
 \label{eq:rs-filtracion}
\end{equation}
\end{proposicion}

\begin{proof}
La sucesión exacta del núcleo y la imagen de cada mapa da
$\length Q_i\le\kappa_i+\sum_jm_{ij}\length S_j$.
Para cada error existe una sobreyección
$(\mathcal O/\varpi^{a_h}\mathcal O)^{b_h}\twoheadrightarrow E_h$,
por lo que $\length E_h\le a_hb_h$.
La aditividad de longitudes a lo largo de la filtración suma
estas desigualdades.
\end{proof}

A93, p.~63, formula condicionalmente una comparación por graduados
con inyecciones en varios módulos y errores aniquilados por una
potencia fija y con un número acotado de generadores.
La Proposición~\ref{prop:rs-filtracion} incluye además la variante
de núcleos controlados $\kappa_i$. B76, p.~54, conserva las
multiplicidades $m_\pi=\dim_{K_\lambda}V_\pi$ en la cota
del denominador; en p.~55
sitúa el retículo $T_\pi$ en
$\rho_\lambda\otimes\xi_1^{-1}$ y distingue la condición
finita de Selmer. Los grupos de torcido negativo
aportados por las líneas de caracteres, con condición local cero,
se definen en p.~56 y contribuyen junto al término asociado a
$A_\pi$; el balance de índices y sus normalizaciones continúa
en p.~57.

La aplicación theta exige identificar $D$, construir
su filtración y los mapas \eqref{eq:rs-graduados}, y acotar los
$\kappa_i,a_h,b_h,m_{ij}$ con la uniformidad requerida.
Si $e_\lambda=v_\varpi(\ell)$, el paso de longitud sobre
$\mathcal O$ a la valoración normalizada usa
$\length_{\mathcal O}/e_\lambda$, como en las notas anteriores.
La finitud de cada módulo individual tampoco proporciona por
sí sola una cota uniforme para una familia variable.

\subsection{Lo demostrado y el siguiente paso}

El Teorema~\ref{thm:rs-inyeccion} construye la clase de anulador
exacto $(\varpi^n)$ y prueba su supervivencia divisible;
la Proposición~\ref{prop:rs-local} calcula un núcleo mediante
datos locales. El ejemplo con $37$ realiza una clase estricta
no nula, con la existencia residual y la escisión local de Ribet
como entradas externas.

La comparación filtrada de la familia theta sigue pendiente,
con sus multiplicidades, torcidos y errores.
La Nota~\ref{sec:finitud-uniforme} retoma las longitudes de la
Nota~\ref{sec:cohomologia}: prueba la finitud para datos fijos
mediante teoría $K$ y desarrolla correcciones locales y cotas
de torre. Ese estudio no sustituye la construcción de los
mapas de la comparación theta.

\subsection{Ejercicios resueltos}

\begin{ejercicio}\label{ej:rs-nucleo}
Sea $\eta\ne1$ y suponga que el ideal generado por $\eta(g)-1$
es $(\varpi^2)$. Calcule el núcleo del cambio de coeficientes
para $n=4$. Compárelo con el caso $\eta=1$ y explique por qué
la igualdad de caracteres residuales no decide por sí sola
la inyectividad.
\end{ejercicio}

\paragraph{Solución del Ejercicio~\ref{ej:rs-nucleo}.}
En el primer caso $A^G=A[\varpi^2]$ y $\varpi^4A^G=0$.
El núcleo es isomorfo a $\mathcal O/\varpi^2\mathcal O$ y
tiene longitud $2$. Si $\eta=1$, en cambio, $A^G=A$ es
divisible y $A^G/\varpi^4A^G=0$; el mapa es inyectivo.
Ambos caracteres son residualmente triviales, pero sus módulos
de invariantes tienen comportamientos distintos bajo multiplicación
por $\varpi^4$.

\begin{ejercicio}\label{ej:rs-local}
En la Proposición~\ref{prop:rs-local}, tome $n=5$ y dos lugares
con segundas coordenadas de valoración $1$ y $3$.
Describa el mapa construido, su núcleo y su imagen.
\end{ejercicio}

\paragraph{Solución del Ejercicio~\ref{ej:rs-local}.}
Se toma $t=3$ y el mapa es $u\mapsto u\varpi^3\gamma_5$.
Su núcleo es $\varpi^2\mathcal O/\varpi^5\mathcal O$, de
longitud $3$, y la imagen tiene longitud $2$.
La suma $1+3$ no es la pérdida necesaria: una misma multiplicación
por $\varpi^3$ anula ambas restricciones. No se deduce que el
Selmer completo tenga longitud $2$.

\begin{ejercicio}\label{ej:rs-tres}
Para $N>n\ge1$, considere
\[
 \mathbb T=\{(a,b,c)\in\mathcal O^3:
               a\equiv b\pmod{\varpi^n},\
               a\equiv c\pmod{\varpi^N}\}.
\]
Calcule $I_{1,2}$ y $\mathfrak c_1$ de
\eqref{eq:rs-ideales-hecke}. Explique por qué una clase de
longitud $n$ no basta para controlar el proyector sobre la
primera coordenada.
\end{ejercicio}

\paragraph{Solución del Ejercicio~\ref{ej:rs-tres}.}
Las diferencias $a-b$ generan $(\varpi^n)$: están en ese ideal
y $(\varpi^n,0,\varpi^n)\in\mathbb T$ lo alcanza.
Para que $(a,0,0)\in\mathbb T$ hacen falta simultáneamente
$a\in(\varpi^n)$ y $a\in(\varpi^N)$; así
$\mathfrak c_1=(\varpi^N)$. El cociente de proyección es
$\mathcal O/\varpi^N\mathcal O$, pues la primera proyección de
$\mathbb T$ es todo $\mathcal O$. Su longitud $N$ puede ser
mucho mayor que $n$.

\begin{ejercicio}\label{ej:rs-signo}
Para el ejemplo $p=37$, $k=32$, determine el torcido del cociclo,
su carácter residual y el valor residual en Frobenius en $2$.
Explique qué parte de la no nulidad se recibe de Ribet y cuál
se demuestra mediante el cambio de coeficientes.
\end{ejercicio}

\paragraph{Solución del Ejercicio~\ref{ej:rs-signo}.}
El cociclo pertenece a $\operatorname{Hom}(\omega^{31},1)$;
el torcido es $-31$, con carácter
$\omega^{-31}=\omega^5$. En Frobenius en $2$ vale
$2^5=32\pmod{37}$. Ribet proporciona la representación residual
no escindida y su escisión local en $37$.
El Lema~\ref{lem:rs-coeficientes} prueba que su coordenada no nula
sobre $\mathbf F_{37}$ sobrevive en
$\Q_{37}/\Z_{37}(-31)$; \eqref{eq:rs-local-q} completa la
trivialidad en los otros primos.

\begin{ejercicio}\label{ej:rs-multiplicidades}
Suponga que una filtración tiene un graduado que se inyecta en
$S^{\oplus3}$, otro que se inyecta en $S'$ y un error anulado
por $\varpi^2$ y generado por cuatro elementos. Si
$\length S=2$ y $\length S'=5$, dé una cota para la longitud
total. ¿Qué dato faltaría para hacerla uniforme en una familia?
\end{ejercicio}

\paragraph{Solución del Ejercicio~\ref{ej:rs-multiplicidades}.}
La cota es $3\cdot2+5+2\cdot4=19$.
Para una familia hay que controlar conjuntamente las longitudes
de $S,S'$, las multiplicidades y los exponentes y números de
generadores de los errores; si los mapas dejan de ser inyectivos,
también sus núcleos. La existencia individual de una filtración
finita no da esas cotas uniformes.
